% ================================================================
%  conteudo/material_metodos.tex
% ================================================================

\section{Material e Métodos}

Esta seção deve descrever, de forma clara e detalhada, todos os procedimentos experimentais realizados durante a prática, de modo que outro pesquisador seja capaz de reproduzi-los integralmente. O texto deve ser redigido no pretérito e na voz passiva, sem uso de verbos na primeira pessoa. Abaixo é apresentado um exemplo de descrição experimental completa.

\subsection{Microrganismo e preparo do inóculo}

Utilizou-se a levedura \textit{Saccharomyces cerevisiae} (linhagem ATCC 2602) como microrganismo modelo. O inóculo foi preparado a partir de cultura estoque mantida em meio YPD (extrato de levedura 10~g~L$^{-1}$, peptona 20~g~L$^{-1}$ e glicose 20~g~L$^{-1}$), incubada a 30~°C por 18~h sob agitação de 150~rpm. A concentração celular inicial no reator foi ajustada para 0,5~g~L$^{-1}$ de biomassa seca por meio de centrifugação e ressuspensão.

\subsection{Condução do processo fermentativo}

\subsubsection{Preparo do meio de cultivo}

O meio de cultivo sintético utilizado foi composto por glicose (fonte de carbono), sulfato de amônio (fonte de nitrogênio) e sais minerais. Foram testadas três concentrações iniciais de glicose: 10, 20 e 40~g~L$^{-1}$. O pH do meio foi ajustado para 5,5 com solução de HCl (1~mol~L$^{-1}$) antes da esterilização em autoclave a 121~°C por 20~min.

\paragraph{Esterilização e montagem do reator}

O reator de bancada (volume útil de 2~L) foi esterilizado separadamente e conectado ao sistema de aeração, controle de temperatura e agitação mecânica antes da inoculação. Todo o procedimento de montagem foi realizado em condições assépticas sob câmara de fluxo laminar, conforme descrito por Schmidell \textit{et al.} (2001).

\subsection{Amostragem e métodos analíticos}

Amostras de 5~mL foram coletadas a cada 2~h durante 24~h de fermentação. Para a determinação da concentração de biomassa, mediu-se a absorbância a 600~nm em espectrofotômetro (modelo UV-1800, Shimadzu) e os valores foram convertidos em concentração de massa seca (g~L$^{-1}$) por meio de curva de calibração previamente estabelecida. A concentração de glicose residual foi determinada pelo método colorimétrico do DNS, conforme Miller (1959). Cada análise foi realizada em triplicata e os resultados expressos como média $\pm$ desvio padrão.

